Nečistoty a poruchy mřížky rozptylují elektrony, snižují pohyblivost a vytvářejí \textbf{zbytkovou rezistivitu $\rho_0$}:
\[\rho(T)\approx\rho_0+\rho_{\mathrm{fonony}}(T).\]
Při nízké teplotě slábne rozptyl na tepelných kmitech, takže převládne $\rho_0$. Čím více nečistot, tím vyšší nenulové plató rezistivity při $T\to0$; čistší kov má nižší $\rho_0$.
